\specialchapter{自伴性判据}

这里汇总正文中陈述的、可用来证明某些部分算子自伴的判据。

固定一个复希尔伯特空间 E 以及 E 上的部分算子 u。在下列每一种情形中，该算子都是自伴的：

\begin{enumerate}
\def\labelenumi{(\roman{enumi})}
\item
  存在一个单射且为埃尔米特的 \(v \in \mathcal{L}(\mathrm{E})\) ，使得 \(u = v^{-1}\) 。（IV，第239页命题9）；
\item
  算子 \(u\) 是对称的，且 \(\mathrm{dom}(u) = \mathrm{E}\) （IV，第240页命题10）；
\item
  算子 \(u\) 是对称的，并诱导从 \(\mathrm{dom}(u)\) 到 E 的双射（IV，第240页命题10的推论）；
\item
  存在 \(\lambda \in \mathbf{C}\) ，使得 \(u + \lambda 1_{\mathrm{E}}\) 和 \(u + \overline{\lambda} 1_{\mathrm{E}}\) 均为满射（IV，第240页命题11）；
\item
  存在一个定义域稠密的闭部分算子 \(v\) ，使得 \(u = v^{*} \circ v\) （IV，第241页命题12）；
\item
  算子 \(u\) 是对称且闭的，并且其亏指标为 \((0,0)\) （IV，第257页命题26的推论），也就是说，\(\mathrm{Ker}(u^{*} - i) = \mathrm{Ker}(u^{*} + i) = \{0\}\) ；
\item
  空间 E 具有 \(\mathrm{L}^{2}(\mathrm{X},\mu)\) 的形式，其中 X 是局部紧拓扑空间，\(\mu\) 是 X 上的测度，而 u 是乘法算子，即乘以一个在 X 上 \(\mu\) -可测且局部地 \(\mu\) -几乎处处取实值的函数（命题23的推论）；
\item
  存在 E 上的正规部分算子 v，以及一个取实值的普遍可测函数 \(f \in \mathcal{L}_{\mathrm{u}}(\mathrm{Sp}(v))\) ，使得 \(u = f(v)\) （IV，第273页推论1）；
\item
  存在 E 上的部分埃尔米特形式 q，使 u 是表示 q 的部分算子（IV，第293页命题14）；
\item
  有 \(u = v + w\) ，其中 v 是自伴部分算子，w 是关于 v 相对有界且满足 \(\|w\|_{v} < 1\) 的对称部分算子（IV，第306页定理5）；
\item
  算子 \(u\) 是 E 上的酉表示 \(\varrho\) （其群为 \(\mathbf{R}\)）的无穷小生成元（V，第428页定理1）。
\end{enumerate}
